
\subsubsection{5.1 Why Baseline Matches Transformer}

The parity between per-layer statistical aggregation and transformer-based token embeddings (both 80\% $\pm$ 4.2\% test accuracy) has three potential explanations.

\textbf{Dataset size effect.} With only 70 training samples, the task may not require deep feature learning. Architecture family classification is coarse-grained---distinguishing ResNet from ViT based on the presence or absence of attention layers, not subtle variations within families. Per-layer statistics (mean, standard deviation, norm) capture these macro-level differences. ResNet exhibits bimodal distributions from batch normalization shift and scale parameters. ViT attention weight matrices exhibit low-rank structure with sparse patterns. EfficientNet depthwise separable convolutions produce factorized kernel statistics. ConvNeXt layer-scaled residuals show controlled weight norms. These family signatures are detectable via shallow aggregation. Larger datasets (500+ models) may reveal within-family variations requiring cross-layer reasoning.

\textbf{Task complexity.} Architecture family classification may favor local patterns over global structure. Harder tasks such as property prediction (test accuracy regression) may require modeling cross-layer interactions (residual flow, gradient propagation). Backdoor detection may need global anomaly detection (backdoor triggers spanning multiple layers). Model merging may require understanding layer-to-layer dependencies (mergeable versus critical layers). The present validation establishes tokenization baseline performance; harder tasks could reveal transformer advantages.

\textbf{Family-level separability.} Strong discriminability at the layer level reduces the need for global reasoning. Per-family accuracy shows that ViT is most distinct (attention mechanism provides clear signal, 85\% transformer accuracy), ResNet and ConvNeXt overlap (both use residual connections, 80\% baseline and 80\% transformer), and EfficientNet is intermediate (shares convolutions with ResNet but has unique scaling, 80\% baseline and 75\% transformer). Layer-wise patterns dominate family discrimination. Cross-layer attention may help only when families share similar layer-level statistics but differ in architectural composition (e.g., block ordering, connection patterns).

\subsubsection{5.2 Implications for Weight-Space Learning}

\textbf{Tokenization is viable.} 80\% $\pm$ 4.2\% test accuracy (versus 60\% threshold and 25\% random baseline) confirms that layer-wise processing preserves sufficient structural signal for property prediction. Future work can build on this foundation with confidence that tokenization does not lose critical information.

\textbf{Baselines matter.} Simple statistical features provide strong comparison points. Researchers proposing novel weight-processing architectures should establish per-layer statistics baselines to isolate gains from architectural inductive biases versus feature quality.

\textbf{Capacity calibration needed.} Transformer overfitting (100\% train, 73.33\% validation) highlights model-data mismatch. Weight-space learning datasets are small compared to vision or NLP corpora. Architecture choices must account for limited sample sizes. For small datasets (100 models), per-layer statistics or shallow transformers are appropriate. For medium datasets (500--1000 models), regularized transformers with dropout and weight decay are suitable. For large datasets (5000+ models), full-capacity transformers with cross-layer attention may be justified.

\subsubsection{5.3 Limitations}

\textbf{Dataset limitations.} 100 models (70 train) limits conclusions about transformer benefits. Cross-layer attention advantages may emerge on larger datasets. All models are image classifiers (ResNet, ViT, EfficientNet, ConvNeXt). Generalization to NLP transformers, reinforcement learning policies, or diffusion models is unknown. All models are pretrained on ImageNet. Weight patterns may differ for models trained on medical images, satellite imagery, or other domains.

\textbf{Methodological limitations.} The 60\% threshold was set without a pilot study. Higher thresholds (e.g., 70--80\%) would increase rigor but risk rejecting valid tokenization strategies. Architecture family classification is coarse-grained. Validation on property prediction, backdoor detection, or model merging tasks would strengthen claims. The work classifies ResNet versus ViT but does not distinguish ResNet-50 versus ResNet-101 or ViT-Base versus ViT-Large. Finer-grained tasks may require cross-layer reasoning. The $10\times$ parameter difference (200K versus 2M) confounds comparison of tokenization strategies (statistical aggregation versus learned embeddings) with model capacity effects. A capacity-matched transformer (200K parameters) would better isolate tokenization quality. All experiments use seed 42. Train/test split variance is not estimated. Confidence intervals ($\pm 4.2$\%) reflect bootstrap resampling variance only, not split sensitivity.

\textbf{Interpretation cautions.} Baseline-transformer parity may reflect dataset properties (family separability) rather than fundamental limits of cross-layer attention. Different datasets may show different patterns. Transformer's 100\% training accuracy suggests memorization. With more data, the gap between baseline and transformer may widen. The tokenization strategy uses flatten, pad, and per-layer z-score normalization. Alternative strategies (graph-based tokenization, hierarchical encodings, learned embeddings) may alter the baseline-transformer comparison.

\subsubsection{5.4 Negative Result: Equivariant GNN Mechanism}

An exploratory experiment tested whether equivariant graph neural networks (GNNs) capture local permutation-symmetric patterns complementary to transformer global attention. The hypothesis predicted that GNNs would show greater than 30\% symmetry differential (accuracy drop when neuron indices are shuffled across layers versus within layers), while transformers would show less than 10\% differential due to global permutation invariance.

\textbf{Setup.} 30 timm models were tokenized and classified into 4 bins based on parameter count. A 2-layer transformer (128 hidden dimensions, 4 attention heads) and a simplified 2-layer GNN (linear layers, no true equivariance guarantees) were trained for 10 epochs. Test set accuracy was measured under three conditions: unperturbed, within-layer neuron permutation, and across-layer neuron permutation. Symmetry differential was computed as the difference between within-layer and across-layer perturbed accuracy.

\textbf{Results.} The transformer achieved 60\% accuracy on all three conditions (unperturbed, within-layer, across-layer), yielding 0\% differential. This confirms global permutation invariance. The GNN achieved 80\% unperturbed accuracy, 40\% within-layer accuracy, and 20\% across-layer accuracy, yielding 20\% differential. The transformer passed the gate (differential less than 10\%), but the GNN failed (20\% less than the 30\% threshold).

\textbf{Analysis.} The GNN showed directional evidence of local bias (20\% differential versus transformer's 0\%), but the magnitude was below the threshold. The proof-of-concept implementation used simplified linear layers instead of true E(n)-equivariant GNN layers from torch\_geometric, which lack the permutation equivariance guarantees necessary for stronger local sensitivity. Additional limitations include small dataset size (30 models versus 100+ recommended), short training (10 epochs versus 50 recommended), and simulated labels (parameter count bins versus real test accuracy). The gate failure indicates that either the GNN mechanism requires proper equivariant implementation and larger-scale validation, or the hypothesis that local permutation sensitivity exceeds 30\% may be overestimated for weight-space tasks.

\textbf{Reflection.} The hypothesis that weight-processing backbones capture orthogonal structural properties (global dependencies versus local permutation symmetries) was only partially supported. The transformer global behavior was confirmed (0\% differential), but the GNN local sensitivity magnitude (20\%) fell short of the prediction (30\%). Whether this reflects implementation shortcuts (simplified GNN instead of true EGNN) or a fundamental property of weight-space learning remains unresolved. Consequently, the complementarity hypothesis---that combining transformer and GNN embeddings improves downstream tasks---could not be tested.
